\documentclass[12pt,a4paper]{article}
\usepackage[utf8]{inputenc}
\usepackage[T1]{fontenc}
\usepackage[portuguese]{babel}
\usepackage{geometry}
\usepackage{enumitem}
\usepackage{longtable}
\usepackage{array}
\usepackage{booktabs}
\usepackage{lmodern}
\geometry{margin=2.4cm}
\setlist[itemize]{leftmargin=1.5em}
\setlist[enumerate]{leftmargin=1.8em}
\title{Engenharia de Software\\Sprint 1 --- Análise do Produto}
\author{}
\date{28/09 \quad | \quad Sprint: 1 \quad | \quad Projeto: Aplicação Universitária UBI}
\begin{document}
\maketitle
\section{Visão e âmbito}
\textbf{Problema:} a informação e os serviços académicos encontram-se dispersos por diferentes canais, dificultando tarefas quotidianas e a orientação no campus.\par
\textbf{Proposta:} centralizar o acesso à informação e aos serviços universitários numa aplicação, considerando estudantes, docentes e funcionários.\par
\textbf{Âmbito inicial a validar:} consulta de horários e salas, notificações de alterações e comunicação docente--estudante. As restantes funcionalidades podem ser consideradas para versões futuras, conforme prioridades, dependências e feedback dos utilizadores.
\section{Stakeholders}
\begin{itemize}
\item \textbf{Estudantes} Utilizadores principais; consultam horários, notas e a cantina.
\item \textbf{Docentes} Gerem aulas, avaliações e comunicação com as turmas.
\item \textbf{Funcionários administrativos} Gerem inscrições, propinas e processos administrativos.
\item \textbf{Serviços Académicos (SASUBI/Reitoria)} Responsáveis pelos processos.
\item \textbf{Equipa de TI da UBI} Manutenção, suporte técnico e integrações.
\item \textbf{Serviços de Ação Social} Gerem bolsas, refeitório e residências.
\item \textbf{Segurança/Portaria} Controlo de acessos aos edifícios, caso aplicável.
\item \textbf{Fornecedores externos} Prestadores potencialmente integrados, como cantina e transportes.
\item \textbf{Visitantes/Candidatos} Acesso limitado à informação pública e às candidaturas.
\end{itemize}
\section{Funcionalidades identificadas}
\begin{enumerate}
\item Consulta de horários e salas.
\item Consulta de notas.
\item Consulta de propinas.
\item Notificações sobre alterações de aulas e frequências.
\item Reservas de espaços.
\item Mapa interativo da universidade, com salas e serviços da UBI.
\item Comunicação entre docentes e estudantes.
\item Gestão de pedidos de documentos.
\end{enumerate}
\section{Personas}
\subsection{Inês Carvalho, estudante}
19 anos; frequenta o 1.º ano de Engenharia Informática. É de Viseu, estuda na Covilhã e vive numa residência. Usa o telemóvel para quase tudo.\par
\textbf{Objetivos e frustrações:} quer consultar o horário, as salas e as notas num só local, ver o menu da cantina e orientar-se pelos vários polos. Frustra-se por alternar entre sites e não saber a quem recorrer para assuntos administrativos.
\subsection{Professora Helena Duarte, docente}
52 anos; docente de Matemática há mais de 20 anos. Tem conhecimentos tecnológicos medianos e prefere ferramentas simples.\par
\textbf{Objetivos e frustrações:} quer comunicar rapidamente com os alunos, lançar sumários e notas, e reservar salas e equipamentos. Frustra-se com tarefas repetitivas e com a utilização de sites diferentes.
\subsection{Paula Rodrigues, funcionária dos Serviços Académicos}
45 anos; atende estudantes presencialmente e por email. O volume de trabalho aumenta no início do semestre e nas inscrições para exames.\par
\textbf{Objetivos e frustrações:} quer reduzir pedidos repetitivos e organizar os processos digitalmente. Frustra-se com filas, pedidos incompletos e perguntas cuja resposta já está publicada.
\subsection{João Almeida, estudante com mobilidade reduzida}
20 anos; estudante de Marketing que se desloca em cadeira de rodas.\par
\textbf{Objetivos e frustrações:} quer consultar percursos e edifícios acessíveis, reservar apoio e receber avisos sobre elevadores avariados. Frustra-se quando descobre no local que um acesso não está disponível.
\section{Cenários de utilização}
\subsection{Cenário 1, Consultar horário e chegar à aula}
Inês acorda de manhã sem saber se houve alterações ao horário. Abre a aplicação e consulta as aulas do dia, incluindo horas e salas. Antes de sair de casa, verifica se há uma notificação de alteração. Ao chegar à universidade, consulta a localização da sala e segue para a aula. Se a informação estiver desatualizada ou a sala não estiver identificada, a aplicação deve indicar que os dados não estão disponíveis e apontar um canal de contacto.
\subsection{Cenário 2, Comunicar uma alteração de aula}
A professora Helena precisa de informar a turma de que a aula mudou de sala. Acede à funcionalidade de comunicação, seleciona a turma e a aula, escreve a alteração e confirma o envio. Os estudantes recebem uma notificação com a informação essencial. Helena deve conseguir verificar que a mensagem foi enviada; se não tiver autorização ou a turma não estiver disponível, o sistema deve explicar o problema sem enviar a mensagem.
\subsection{Cenário 3, Pedir um documento académico}
Paula recebe muitos pedidos de declarações, alguns sem os dados necessários. O estudante consulta os documentos disponíveis, seleciona o tipo pretendido e preenche os campos obrigatórios. Antes do envio, revê o pedido. O sistema confirma a submissão e apresenta um identificador ou estado para acompanhamento. Se faltarem dados, assinala os campos a corrigir.
\subsection{Cenário 4, Encontrar um percurso acessível}
João precisa de chegar a um edifício num polo que não conhece. Consulta o mapa, pesquisa o edifício e seleciona a opção de percurso acessível. A aplicação apresenta um percurso que considera a existência de entradas acessíveis. Se um elevador estiver avariado, o aviso deve ser visível e o sistema deve apresentar uma alternativa quando exista; caso contrário, deve informar que não consegue garantir um percurso acessível.
\subsection{Cenário 5, Consultar uma nota ou propina}
Inês quer confirmar uma nota publicada ou verificar informação sobre propinas. Abre a secção correspondente e consulta os dados associados à sua conta. A aplicação distingue informação atual de informação indisponível ou ainda não publicada e não expõe dados de outros estudantes.
\section{User Stories iniciais}
\begin{itemize}
\item \textbf{US01, Horário:} Como estudante, quero consultar o meu horário e as salas, para planear o meu dia sem alternar entre vários sites.
\item \textbf{US02, Notificações:} Como estudante, quero receber avisos de alterações a aulas ou frequências, para agir com base em informação atualizada.
\item \textbf{US03, Comunicação:} Como docente, quero comunicar uma alteração à turma, para que os estudantes sejam informados atempadamente.
\item \textbf{US04, Percurso acessível:} Como estudante com mobilidade reduzida, quero consultar percursos acessíveis e avisos sobre elevadores, para planear a deslocação.
\item \textbf{US05, Pedido de documento:} Como estudante, quero submeter e acompanhar um pedido de documento, para evitar deslocações e saber o estado do pedido.
\item \textbf{US06, Gestão de pedidos:} Como funcionária dos Serviços Académicos, quero receber pedidos com os campos necessários, para reduzir pedidos incompletos.
\end{itemize}
\noindent\textbf{Nota:} estas histórias são candidatas ao Product Backlog. A equipa deve refiná-las, atribuir prioridade e selecionar apenas as adequadas à Sprint, com critérios de aceitação acordados.
\section{Priorização preliminar}
\small
\begin{longtable}{>{\centering\arraybackslash}p{0.7cm} p{3.4cm} p{1.5cm} p{7.8cm}}
\toprule
\textbf{N.º} & \textbf{Funcionalidade} & \textbf{Prioridade} & \textbf{Justificação}\\
\midrule
1 & Horários e salas & Alta & Uso recorrente e informação essencial para estudantes.\\
2 & Notas & Alta & Apoia o acompanhamento do semestre.\\
3 & Propinas & Média & Importante, mas consultada com menor frequência.\\
4 & Notificações de alterações & Alta & Ajuda a agir perante mudanças e evita deslocações ou faltas.\\
5 & Reservas de espaços & Baixa & Útil; prioridade deve ser validada com utilizadores e âmbito.\\
6 & Mapa interativo & Baixa & Pode ser mais valioso para novos alunos, visitantes e pessoas com necessidades de acessibilidade.\\
7 & Comunicação docente--estudante & Alta & Facilita a comunicação direta sobre aulas e atividades.\\
8 & Pedidos de documentos & Média & Pode reduzir trabalho administrativo e pedidos repetitivos.\\
\bottomrule
\end{longtable}
\normalsize
\section{Questões em aberto e validação}
\begin{itemize}
\item Confirmar o âmbito e o nome oficial do produto, os utilizadores-alvo e se a aplicação integra sistemas existentes ou apenas agrega ligações/informação.
\item Validar personas e cenários com estudantes, docentes e funcionários; até essa validação, tratá-los como proto-personas e hipóteses, não como factos observados.
\item Confirmar como serão obtidos horários, notas, propinas, dados de acessibilidade e estado de elevadores, bem como quem mantém esses dados.
\item Rever as prioridades: o valor para utilizadores com necessidades específicas não deve ser confundido com a frequência de utilização pela maioria.
\item Definir critérios de aceitação para as histórias escolhidas e associá-las às funcionalidades; evitar colocar todas as funcionalidades na Sprint 1 sem capacidade e objetivos claros.
\end{itemize}
\end{document}
